\section{The source calculus \texorpdfstring{$\CICi$}{CIC-ix}}
\label{sec:calculus}

This section defines the source calculus: a monomorphic, first-order
fragment of the Calculus of Inductive Constructions whose inductive
datatypes are \emph{indexed families}.  As in the companion paper
\cite{Extraction2026}, the presentation is \emph{stratified}: set types,
propositions, predicate kinds, terms and proofs are separate syntactic
categories, so the outcome of CIC's sort analysis --- which subterms are
computational, which are logical --- is built into the grammar, and the
erasure of \cref{sec:lambdabox} and the translation of
\cref{sec:translation} are definable by structural recursion on raw
syntax, with no typing oracle.  The fragment of \cite{Extraction2026} is
literally the sub-calculus in which every datatype has zero indices and
premise-free constructors with closed argument types;
\cref{sec:calculus-cic} discusses how the extended fragment is carved out
of full CIC.

\subsection{Syntax}
\label{sec:calculus-syntax}

We fix countably infinite, pairwise disjoint sets of \emph{term variables}
(written $x, y, z, w, f, \dots$), of \emph{proof variables} ($p, q, h,
\dots$) and of \emph{predicate variables} ($X, \dots$), and disjoint sets
of names for \emph{datatypes} ($I$), for \emph{datatype constructors}
($\wcns{c}, \wcns{d}$), for \emph{inductive predicates} ($J$), for their
\emph{clause constructors} ($\wcns{k}$), and for \emph{constants} ($c$).

\begin{definition}[Syntax of $\CICi$]
\label{def:syntax}
The syntactic categories of \emph{set types} $\sigma, \tau$,
\emph{propositions} $\varphi, \psi$, \emph{predicate kinds} $K$,
\emph{predicate expressions} $P$, \emph{terms} $t, a, b, u, v$ and
\emph{proofs} $\pi$ are defined by the following grammar.
\begin{align*}
  \sigma, \tau &\Coloneqq
    I(\tup{a})
    \mid \Pi x{:}\sigma.\,\tau
    \mid \varphi \to \tau
    \mid \refty{x{:}\sigma}{\varphi}
  \\
  \varphi, \psi &\Coloneqq
    \top \mid \bot
    \mid \varphi \wedge \psi
    \mid \varphi \vee \psi
    \mid \varphi \to \psi
    \mid \forall x{:}\sigma.\,\varphi
    \mid \exists x{:}\sigma.\,\varphi
  \\ &\mathrel{\phantom{\Coloneqq}}
    \mid \forall X{:}K.\,\varphi
    \mid a =_{\sigma} b
    \mid J(\tup{a})
    \mid X(\tup{a})
  \\
  K &\Coloneqq (\sigma_1, \dots, \sigma_n) \to \Prop \qquad (n \geq 1)
  \\
  P &\Coloneqq X \mid \lambda(x_1{:}\sigma_1, \dots, x_n{:}\sigma_n).\,\varphi
  \\
  t, a, b &\Coloneqq
    x
    \mid c
    \mid t\,a
    \mid t\,\pi
    \mid \lambda x{:}\sigma.\,t
    \mid \lambda p{:}\varphi.\,t
    \mid \wcns{c}(\tup{a}; \tup{\pi})
  \\ &\mathrel{\phantom{\Coloneqq}}
    \mid \tcase{a}{\tup{z}.z.\tau}{\tup{br}}
    \mid \tfix\,f{:}T.\,t
    \mid \trefine{t}{\pi}
    \mid \tval{t}
  \\ &\mathrel{\phantom{\Coloneqq}}
    \mid \tsplit{\pi}{p_1\,p_2.\,t}
    \mid \tcast{\pi}{z.\tau}{t}
    \mid \tscase{J}{\pi}{\tup{z}.\tau}{\tup{p}.\,t}
    \mid \texfalso{\pi}{\sigma}
  \\
  \pi &\Coloneqq
    p
    \mid c
    \mid \lambda p{:}\varphi.\,\pi
    \mid \pi\,\pi'
    \mid \lambda x{:}\sigma.\,\pi
    \mid \pi\,a
    \mid \Lambda X{:}K.\,\pi
    \mid \pi\,P
    \mid \langle \pi_1, \pi_2 \rangle
    \mid \pi.1
    \mid \pi.2
  \\ &\mathrel{\phantom{\Coloneqq}}
    \mid \iota_1(\pi) \mid \iota_2(\pi)
    \mid \mathsf{vcase}(\pi;\, p_1.\pi_1;\, p_2.\pi_2)
    \mid \langle a, \pi \rangle_{\exists x{:}\sigma.\varphi}
    \mid \mathsf{ecase}(\pi;\, x\,p.\,\pi')
  \\ &\mathrel{\phantom{\Coloneqq}}
    \mid \mathsf{tt}
    \mid \texfalso{\pi}{\varphi}
    \mid \mathsf{refl}_{\sigma}(a)
    \mid \mathsf{rew}(\pi;\, z.\psi;\, \pi')
    \mid \mathsf{discr}(\pi)
    \mid \wcns{k}(\tup{a}; \tup{\pi})
  \\ &\mathrel{\phantom{\Coloneqq}}
    \mid \tind{J}(\tup{z}.\psi;\, \tup{\pi};\, \pi)
    \mid \tind{I}(\tup{z}.z.\psi;\, \tup{\Pi};\, a)
    \mid \tvalid{t}
  \end{align*}
where:
\begin{itemize}[leftmargin=2em]
\item a branch list $\tup{br}$ has the form
  $\{\wcns{c}_1(\tup{y}_1; \tup{p}_1; \tup{q}_1) \Rightarrow t_1 \mid
  \dots \mid \wcns{c}_M(\tup{y}_M; \tup{p}_M; \tup{q}_M) \Rightarrow
  t_M\}$ with one branch for every constructor of the datatype concerned;
  in the branch for $\wcns{c}_i$, the term variables $\tup{y}_i$ bind the
  constructor's informative arguments, the proof variables $\tup{p}_i$
  its propositional premises, and the proof variables $\tup{q}_i$ (one
  per index of the datatype) the \emph{index equations} of the ford;
\item in $\tcase{a}{\tup{z}.z.\tau}{\tup{br}}$ the term variables
  $\tup{z}$ (one per index) and $z$ bind into the motive $\tau$; in
  $\tscase{J}{\pi}{\tup{z}.\tau}{\tup{p}.t}$ the variables $\tup{z}$ (one
  per index of $J$) bind into $\tau$ and the proof variables $\tup{p}$
  bind into $t$; in $\tind{I}$ each clause proof $\Pi_i$ has the form
  $\tup{y}_i\,\tup{p}_i\,\tup{h}_i.\,\pi_i$, binding the constructor
  telescope, its premises, and one induction hypothesis per recursive
  argument;
\item $T$ in $\tfix\,f{:}T.\,t$ ranges over set types;
  $\refty{x{:}\sigma}{\varphi}$ binds $x$ in $\varphi$; the binding
  structure of the remaining forms is the evident one.
\end{itemize}
Terms of the form $\trefine{t}{\pi}$ carry a refinement-type annotation,
and $\langle a, \pi\rangle$ carries its $\exists$-type; we omit the
annotations when irrelevant.  We write $\neg\varphi$ for $\varphi \to
\bot$ and $\varphi \leftrightarrow \psi$ for $(\varphi \to \psi) \wedge
(\psi \to \varphi)$.  All categories are identified up to
$\alpha$-equivalence; substitution is capture-avoiding.
\end{definition}

\begin{notation}
\label{not:vectors}
For an $m$-indexed datatype or predicate we abbreviate index tuples as
$\tup{u} = u_1, \dots, u_m$, and write $\tup{u} =_{\tup{\sigma}} \tup{v}$
for the list of propositions $u_1 =_{\sigma_1} v_1, \dots, u_m
=_{\sigma_m} v_m$; a proof-variable telescope $\tup{q} : (\tup{u}
=_{\tup{\sigma}} \tup{v})$ declares $q_h : u_h =_{\sigma_h} v_h$ for $1
\le h \le m$.  When a constructor has no premises and the datatype no
indices we write $\wcns{c}(\tup{a})$ for $\wcns{c}(\tup{a};)$ and
$\wcns{c}(\tup{y}) \Rightarrow t$ for a branch, recovering the notation
of \cite{Extraction2026}.
\end{notation}

Compared with the companion calculus, the new material is: datatype
occurrences $I(\tup{a})$ applied to \emph{index terms}; constructor terms
$\wcns{c}(\tup{a}; \tup{\pi})$ carrying proofs of the constructor's
propositional premises; the \emph{forded} case former, whose branches
bind index-equation proofs $\tup{q}_i$; the singleton elimination
$\tscase{J}{}{}{}$ for index-determined propositional singletons
(\cref{def:forced-singleton}); the discrimination proof former
$\mathsf{discr}$; and the generalized datatype induction $\tind{I}$.  The
three function spaces, the refinement formers, the propositional
singleton eliminations $\tsplit{}{}\!$, $\tcast{}{}{}\!$,
$\texfalso{}{}\!$ and the proof-level logic are unchanged.  Dependency of
types on terms now enters through \emph{three} channels: refinement
payloads, propositions embedded in types, and index positions of datatype
occurrences.  There are still no type-level functions, no universes and
no large elimination, so set types remain \emph{rigid}
(\cref{lem:rigid}).

\subsection{Environments}
\label{sec:calculus-env}

\begin{definition}[Declarations and environments]
\label{def:env}
A \emph{global environment} $\Env$ is a finite list of declarations of
the following three kinds.
\begin{enumerate}[leftmargin=2em]
\item A \emph{datatype declaration}
  \[
    \mathsf{data}\ I : (\sigma_1, \dots, \sigma_m) \to \Set \defeq
    \{\, \wcns{c}_1 : K_1;\ \dots;\ \wcns{c}_M : K_M \,\}
  \]
  with $m \geq 0$ \emph{indices} and $M \geq 1$ constructors, where the
  index types $\sigma_h$ are closed set types over the preceding
  environment, and each \emph{constructor scheme} $K_i$ has the form
  \[
    K_i \;=\; \forall (y_{1}{:}A_{i1}) \cdots (y_{r_i}{:}A_{ir_i}).\
      \varphi_{i1} \Rightarrow \dots \Rightarrow \varphi_{i\ell_i}
      \Rightarrow I(t_{i1}, \dots, t_{im})
  \]
  subject to: each \emph{argument type} $A_{ij}$ is either a
  \emph{recursive occurrence} $I(\tup{v})$ with $\tup{v}$ terms over the
  preceding environment and $y_1, \dots, y_{j-1}$ (of the index types,
  in the sense of \cref{fig:wf}), or a set type over the preceding
  environment and $y_1, \dots, y_{j-1}$ in which $I$ does not occur;
  each \emph{premise} $\varphi_{il}$ is a proposition over the preceding
  environment and $\tup{y}_i$ in which $I$ does not occur; and each
  \emph{index term} $t_{ih}$ is a term over the preceding environment
  and $\tup{y}_i$ with $t_{ih} : \sigma_h$.  We call $\tup{t}_i =
  t_{i1}, \dots, t_{im}$ the \emph{index pattern} of $\wcns{c}_i$ and
  say the argument position $j$ is \emph{recursive} if $A_{ij}$ is a
  recursive occurrence.
\item An \emph{inductive-predicate declaration}
  $\mathsf{pred}\ J : (\tup{\sigma}) \to \Prop$ with a finite list of
  clauses
  \[
    \wcns{k}_j : \forall \tup{x}_j{:}\tup{\tau}_j.\;
      \varphi_{j1} \to \dots \to \varphi_{j\ell_j} \to J(\tup{u}_j)
  \]
  where the index types $\tup{\sigma}$ are closed set types over the
  preceding environment, the $\tup{\tau}_j$ are set types over the
  preceding environment (possibly mentioning $\tup{x}_j$ bound earlier
  in the same telescope) in which $J$ does not occur, the indices
  $\tup{u}_j$ are terms over $\tup{x}_j$ with $u_{jh} : \sigma_h$, and
  each premise $\varphi_{ji}$ is either of the form $J(\tup{v})$ (a
  \emph{recursive premise}, $\tup{v}$ terms over the telescope of the
  index types) or a proposition over the preceding environment and
  $\tup{x}_j$ in which $J$ does not occur.
\item A \emph{definition}, either a \emph{term definition} $c \defeq t :
  \sigma$ with $\sigma$ a closed set type and $t$ a closed term, or a
  \emph{lemma} $c \defeq \pi : \varphi$ with $\varphi$ a closed
  proposition and $\pi$ a closed proof.
\end{enumerate}
All names declared in $\Env$ are distinct.  An environment in which
every constant has a body is \emph{definitional}; this paper considers
definitional environments (see \cref{rem:axioms} for postulates).
A datatype with $m = 0$ and premise-free constructors with closed
argument types is exactly a datatype of \cite{Extraction2026}; we say
$I$ is a \emph{family} when $m > 0$.
\end{definition}

\begin{example}[Running declarations]
\label{ex:declarations}
Fix $\nat$ with $\csO : \nat$-constructors as in \cite{Extraction2026}
(formally: $\mathsf{data}\ \nat : () \to \Set \defeq \{\csO : \nat;\
\csS : \forall y{:}\nat.\,\nat\}$, where we write $\nat$ for the
$0$-index occurrence $\nat()$), $\bool$ likewise, and a datatype $A$
(the element type, e.g.\ $\nat$).  The running families are:
\begin{align*}
  &\mathsf{data}\ \vect : (\nat) \to \Set \defeq\\
  &\qquad \{\ \csvnil : \vect(\csO);\ \
    \csvcons : \forall (n{:}\nat)(a{:}A)(v{:}\vect(n)).\ \vect(\csS\,n)
    \ \}\\
  &\mathsf{data}\ \Fin : (\nat) \to \Set \defeq\\
  &\qquad \{\ \wcns{fz} : \forall n{:}\nat.\ \Fin(\csS\,n);\ \
    \wcns{fs} : \forall (n{:}\nat)(k{:}\Fin(n)).\ \Fin(\csS\,n)\ \}\\
  &\mathsf{data}\ \mathsf{bnd} : (\nat) \to \Set \defeq
    \{\ \wcns{B} : \forall (n{:}\nat)(k{:}\nat).\
      \mathsf{lt}(k, n) \Rightarrow \mathsf{bnd}(n)\ \}\\
  &\mathsf{data}\ \mathsf{rfl}_{\varphi_0} : (\bool) \to \Set \defeq\\
  &\qquad \{\ \wcns{RT} : \varphi_0 \Rightarrow
    \mathsf{rfl}_{\varphi_0}(\cstrue);\ \
    \wcns{RF} : \neg\varphi_0 \Rightarrow
    \mathsf{rfl}_{\varphi_0}(\csfalse)\ \}\\
  &\mathsf{pred}\ \mathsf{le} : (\nat, \nat) \to \Prop \defeq\\
  &\qquad \{\ \wcns{le_n} : \forall n{:}\nat.\ \mathsf{le}(n,n);\\
  &\qquad\phantom{\{}\ \wcns{le_S} : \forall (n\,m{:}\nat).\
      \mathsf{le}(n,m) \to \mathsf{le}(n, \csS\,m)\ \}\\
  &\mathsf{pred}\ \mathsf{eq}_\sigma : (\sigma, \sigma) \to \Prop \defeq
    \{\ \wcns{refl}_{\mathsf{eq}} : \forall x{:}\sigma.\
      \mathsf{eq}_\sigma(x,x)\ \}
\end{align*}
where $\mathsf{lt}$ is some previously declared inductive predicate and
$\varphi_0$ a closed proposition.  $\vect$ and $\Fin$ are informative
recursive families; $\mathsf{bnd}$ is a subset-shaped family (one
constructor, informative carrier, propositional payload, index);
$\mathsf{rfl}_{\varphi_0}$ is an enumeration-shaped family (constant
constructors, propositional payloads, distinct rigid index patterns) ---
the monomorphic image of ssreflect's $\mathsf{reflect}\ P\ b$ at $P =
\varphi_0$; $\mathsf{le}$ is a predicate family;
$\mathsf{eq}_\sigma$ is a user-declared equality --- the primitive
$=_\sigma$ remains available, and \cref{ex:eq-pred} compares the two.
\end{example}

\begin{remark}[Function-indexed families]
\label{rem:function-index}
The index types $\sigma_h$ of \cref{def:env} are arbitrary closed set
types --- including $\Pi$-types, so families indexed by
\emph{functions} are in the fragment.  This is deliberate and sound:
index equality means conversion-class equality in the model
(\cref{def:interp}), and nothing in the development inspects the shape
of index types.  It deserves a warning, though, because function
indices are exactly where the intensionality of that equality bites
hardest (\cref{rem:funext-leak}): the instances $I(\lambda
x{:}\nat.\,\mathsf{add}\ \csO\ x)$ and $I(\lambda x{:}\nat.\,x)$ have
distinct index classes, their interpretations are disjoint
(\cref{lem:index-reflection}(3)), and no emitted axiom relates them
--- so goals that transport members between extensionally equal but
non-convertible indices are soundly, and irremediably within this
model, out of reach.
\end{remark}

\begin{definition}[Contexts]
A \emph{context} $\Gamma$ is a finite list of declarations $x : \sigma$,
$p : \varphi$ and $X : K$ with pairwise distinct variables, each type,
proposition or kind well formed over the preceding declarations
(\cref{fig:wf}).
\end{definition}

\subsection{Typing}
\label{sec:calculus-typing}

The judgements are: $\vdash \Env \wf$; and, relative to a fixed
well-formed $\Env$ left implicit, $\Gamma \vdash \sigma \styp$, $\Gamma
\vdash \varphi \sprop$, $\Gamma \vdash P : K$, $\Gamma \vdash t :
\sigma$, and $\Gamma \vdash \pi : \varphi$.  The rules are given in
\cref{fig:wf,fig:terms,fig:terms2,fig:proofs,fig:proofs2}.  Environment well-formedness requires
each declaration to be well formed over its prefix: for a datatype
declaration, each index type must satisfy $\vdash \sigma_h \styp$, and
for each constructor scheme, working in the telescope $y_1 : A_{i1},
\dots$ extended left to right, each argument type must satisfy $\Gamma_j
\vdash A_{ij} \styp$ (where for a recursive occurrence $I(\tup{v})$ this
means $\Gamma_j \vdash v_h : \sigma_h$ for all $h$, the formation rule
for $I$-occurrences being available at this point for the declaration
being processed), each premise $\Gamma_{r_i} \vdash \varphi_{il}
\sprop$, and each index term $\Gamma_{r_i} \vdash t_{ih} : \sigma_h$;
predicate declarations are analogous; a definition $c \defeq t : \sigma$
(resp.\ $c \defeq \pi : \varphi$) requires $\vdash \sigma \styp$ and
$\vdash t : \sigma$ (resp.\ $\vdash \varphi \sprop$ and $\vdash \pi :
\varphi$) over the prefix.

\begin{figure}[tp]
\small
\[
\begin{array}{c}
\rul{I : (\tup{\sigma}) \to \Set \in \Env \quad
     \Gamma \vdash \tup{a} : \tup{\sigma}}
    {\Gamma \vdash I(\tup{a}) \styp}
\qquad
\rul{\Gamma \vdash \sigma \styp \quad \Gamma, x{:}\sigma \vdash \tau \styp}
    {\Gamma \vdash \Pi x{:}\sigma.\tau \styp}
\\[2.2ex]
\rul{\Gamma \vdash \varphi \sprop \quad \Gamma \vdash \tau \styp}
    {\Gamma \vdash \varphi \to \tau \styp}
\qquad
\rul{\Gamma \vdash \sigma \styp \quad \Gamma, x{:}\sigma \vdash \varphi \sprop}
    {\Gamma \vdash \refty{x{:}\sigma}{\varphi} \styp}
\\[2.2ex]
\rul{}{\Gamma \vdash \top \sprop}
\qquad
\rul{}{\Gamma \vdash \bot \sprop}
\qquad
\rul{\Gamma \vdash \varphi \sprop \quad \Gamma \vdash \psi \sprop}
    {\Gamma \vdash \varphi \mathbin{\circledast} \psi \sprop}
\;{\scriptstyle \circledast \in \{\wedge,\vee,\to\}}
\\[2.2ex]
\rul{\Gamma \vdash \sigma \styp \quad \Gamma, x{:}\sigma \vdash \varphi \sprop}
    {\Gamma \vdash Qx{:}\sigma.\varphi \sprop}
\;{\scriptstyle Q \in \{\forall,\exists\}}
\qquad
\rul{\Gamma, X{:}K \vdash \varphi \sprop \quad \Gamma \vdash K\ \mathsf{kind}}
    {\Gamma \vdash \forall X{:}K.\varphi \sprop}
\\[2.2ex]
\rul{\Gamma \vdash a : \sigma \quad \Gamma \vdash b : \sigma}
    {\Gamma \vdash a =_{\sigma} b \sprop}
\qquad
\rul{J : (\tup{\sigma}) \to \Prop \in \Env \quad \Gamma \vdash \tup{a} : \tup{\sigma}}
    {\Gamma \vdash J(\tup{a}) \sprop}
\\[2.2ex]
\rul{X : (\tup{\sigma}) \to \Prop \in \Gamma \quad \Gamma \vdash \tup{a} : \tup{\sigma}}
    {\Gamma \vdash X(\tup{a}) \sprop}
\qquad
\rul{\Gamma \vdash \sigma_i \styp \ (1 \le i \le n)}
    {\Gamma \vdash (\sigma_1,\dots,\sigma_n) \to \Prop\ \mathsf{kind}}
\\[2.2ex]
\rul{X : K \in \Gamma}{\Gamma \vdash X : K}
\qquad
\rul{\Gamma, \tup{x}{:}\tup{\sigma} \vdash \varphi \sprop}
    {\Gamma \vdash \lambda(\tup{x}{:}\tup{\sigma}).\varphi : (\tup{\sigma}) \to \Prop}
\end{array}
\]
\caption{Well-formedness of set types, propositions, kinds and predicate
  expressions.  $\Gamma \vdash \tup{a} : \tup{\sigma}$ abbreviates
  $\Gamma \vdash a_h : \sigma_h$ for all $h$; since index types are
  closed, no substitution occurs in index telescopes.}
\label{fig:wf}
\end{figure}

\begin{figure}[tp]
\small
\[
\begin{array}{c}
\rul{x : \sigma \in \Gamma}{\Gamma \vdash x : \sigma}
\qquad
\rul{(c \defeq t : \sigma) \in \Env}{\Gamma \vdash c : \sigma}
\qquad
\rul{\Gamma \vdash t : \Pi x{:}\sigma.\tau \quad \Gamma \vdash a : \sigma}
    {\Gamma \vdash t\,a : \subst{\tau}{a}{x}}
\\[2.2ex]
\rul{\Gamma \vdash t : \varphi \to \tau \quad \Gamma \vdash \pi : \varphi}
    {\Gamma \vdash t\,\pi : \tau}
\qquad
\rul{\Gamma, x{:}\sigma \vdash t : \tau}
    {\Gamma \vdash \lambda x{:}\sigma.t : \Pi x{:}\sigma.\tau}
\qquad
\rul{\Gamma, p{:}\varphi \vdash t : \tau}
    {\Gamma \vdash \lambda p{:}\varphi.t : \varphi \to \tau}
\\[2.2ex]
\rul{\begin{array}{c}
     \wcns{c}_i : \forall \tup{y}{:}\tup{A}.\ \varphi_1 \Rightarrow \dots
       \Rightarrow \varphi_\ell \Rightarrow I(\tup{t}) \ \in I \\
     \Gamma \vdash a_j : \subst{A_{j}}{\tup{a}_{<j}}{\tup{y}_{<j}}
       \ (1 \le j \le r) \\
     \Gamma \vdash \pi_l : \subst{\varphi_l}{\tup{a}}{\tup{y}}
       \ (1 \le l \le \ell)
     \end{array}}
    {\Gamma \vdash \wcns{c}_i(\tup{a}; \tup{\pi}) :
       I(\subst{\tup{t}}{\tup{a}}{\tup{y}})}
\qquad
\rul{\Gamma \vdash t : \refty{x{:}\sigma}{\varphi}}
    {\Gamma \vdash \tval{t} : \sigma}
\\[2.2ex]
\rul{\begin{array}{c}
     \Gamma \vdash a : I(\tup{u}) \qquad
     \Gamma, \tup{z}{:}\tup{\sigma}, z{:}I(\tup{z}) \vdash \tau \styp \\
     \Gamma,\ \tup{y}_i{:}\tup{A}_i,\ \tup{p}_i{:}\tup{\varphi}_i,\
       \tup{q}_i{:}(\tup{u} =_{\tup{\sigma}} \tup{t}_i)
       \ \vdash\ t_i :
       \tau[\tup{t}_i/\tup{z},\, \wcns{c}_i(\tup{y}_i;\tup{p}_i)/z]
       \ \ (\text{all } i)
     \end{array}}
    {\Gamma \vdash \tcase{a}{\tup{z}.z.\tau}{\tup{br}}
      : \tau[\tup{u}/\tup{z},\, a/z]}
\\[2.2ex]
\rul{\Gamma \vdash \refty{x{:}\sigma}{\varphi} \styp \quad
     \Gamma \vdash t : \sigma \quad
     \Gamma \vdash \pi : \subst{\varphi}{t}{x}}
    {\Gamma \vdash \trefine{t}{\pi} : \refty{x{:}\sigma}{\varphi}}
\qquad
\rul{\Gamma \vdash \pi : \bot \quad \Gamma \vdash \sigma \styp}
    {\Gamma \vdash \texfalso{\pi}{\sigma} : \sigma}
\end{array}
\]
\caption{Typing of terms, part I.  In the constructor rule, argument
  types are instantiated left to right and premises at the full
  argument list.  In the case rule, $\tup{br}$ abbreviates the branch
  list $\{\wcns{c}_i(\tup{y}_i;\tup{p}_i;\tup{q}_i) \Rightarrow
  t_i\}_i$, and $\tup{u} =_{\tup{\sigma}} \tup{t}_i$ are the index
  equations of the ford (\cref{not:vectors}); with $m = 0$ and
  premise-free constructors the rule degenerates to the case rule of
  \cite{Extraction2026}.}
\label{fig:terms}
\end{figure}

\begin{figure}[tp]
\small
\[
\begin{array}{c}
\rul{\Gamma \vdash \pi : \varphi_1 \wedge \varphi_2 \quad
     \Gamma, p_1{:}\varphi_1, p_2{:}\varphi_2 \vdash t : \tau}
    {\Gamma \vdash \tsplit{\pi}{p_1\,p_2.t} : \tau}
\qquad
\rul{\Gamma \vdash t : \sigma \quad \sigma \conv \sigma' \quad
     \Gamma \vdash \sigma' \styp}
    {\Gamma \vdash t : \sigma'}
\\[2.2ex]
\rul{\Gamma \vdash \pi : a =_{\sigma} b \quad
     \Gamma, z{:}\sigma \vdash \tau \styp \quad
     \Gamma \vdash t : \subst{\tau}{a}{z}}
    {\Gamma \vdash \tcast{\pi}{z.\tau}{t} : \subst{\tau}{b}{z}}
\\[2.2ex]
\rul{\begin{array}{c}
     J \text{ a forced singleton with clause }
     \wcns{k} : \forall \tup{x}{:}\tup{\tau}.\ \varphi_1 \to \dots \to
       \varphi_\ell \to J(\tup{u}^0)\\
     \text{and solution } \jmath \ \ (\text{\cref{def:forced-singleton}});
     \qquad \theta = [u_{\jmath(i)}/x_i]_i \\
     \Gamma \vdash \pi : J(\tup{u}) \qquad
     \Gamma, \tup{z}{:}\tup{\sigma} \vdash \tau \styp \qquad
     \Gamma, \tup{p}{:}\theta(\tup{\varphi}) \vdash t :
       \tau[\theta(\tup{u}^0)/\tup{z}]
     \end{array}}
    {\Gamma \vdash \tscase{J}{\pi}{\tup{z}.\tau}{\tup{p}.t} :
       \tau[\tup{u}/\tup{z}]}
\\[2.2ex]
\rul{\begin{array}{c}
     T = \Pi x_1{:}\sigma_1 \dots \Pi x_n{:}\sigma_n.\tau \qquad
     \sigma_k = I(\tup{v}) \qquad \Gamma \vdash T \styp \\
     t = \lambda \tup{x}{:}\tup{\sigma}.\,
         \tcase{x_k}{\tup{z}.z.\tau_0}{\tup{br}} \qquad
     \tau = \tau_0[\tup{v}/\tup{z},\, x_k/z] \\
     \Gamma, f{:}T, \tup{x}{:}\tup{\sigma},
       \tup{y}_i{:}\tup{A}_i, \tup{p}_i{:}\tup{\varphi}_i,
       \tup{q}_i{:}(\tup{v} =_{\tup{\sigma}'} \tup{t}_i) \\
       \qquad\qquad \vdash\ b_i : \tau_0[\tup{t}_i/\tup{z},\,
         \wcns{c}_i(\tup{y}_i;\tup{p}_i)/z]
       \ \ (\text{all } i) \\
     \mathrm{grd}_k(f; \tup{y}_i^{\mathrm{rec}}; b_i) \ \ (\text{all } i)
     \end{array}}
    {\Gamma \vdash \tfix\,f{:}T.\,t : T}
\end{array}
\]
\caption{Typing of terms, part II.  In the $\tscase{J}{}{}{}$ rule,
  $J$, its clause data and $\jmath$ are as in
  \cref{def:forced-singleton}.  In the $\tfix$ rule, $\tup{br}$
  abbreviates the branch list with bodies $b_i$; $x_k$ is the
  designated \emph{guard argument}; $\tup{\sigma}'$ are the index types
  of $I$; $\tup{y}_i^{\mathrm{rec}} \subseteq \tup{y}_i$ denotes the
  recursive argument positions of $\wcns{c}_i$; and the \emph{guard
  condition} $\mathrm{grd}_k(f; V; b)$ requires every occurrence of $f$
  in $b$ to lie at the head of an application spine $f\,a_1 \cdots
  a_n$ occurring as a subterm of $b$, with $a_k \in V$ (so $f$ may
  head such a spine but not occur as an argument, of itself or of
  anything else; \cref{def:discipline} makes the reading formal).  We
  further require $f, \tup{x}, \tup{y}_i$ not to be shadowed in $b_i$.}
\label{fig:terms2}
\end{figure}

\begin{figure}[tp]
\small
\[
\begin{array}{c}
\rul{p : \varphi \in \Gamma}{\Gamma \vdash p : \varphi}
\qquad
\rul{(c \defeq \pi : \varphi) \in \Env}{\Gamma \vdash c : \varphi}
\qquad
\rul{\Gamma, p{:}\varphi \vdash \pi : \psi}
    {\Gamma \vdash \lambda p{:}\varphi.\pi : \varphi \to \psi}
\\[2.2ex]
\rul{\Gamma \vdash \pi : \varphi \to \psi \quad \Gamma \vdash \pi' : \varphi}
    {\Gamma \vdash \pi\,\pi' : \psi}
\qquad
\rul{\Gamma, x{:}\sigma \vdash \pi : \varphi}
    {\Gamma \vdash \lambda x{:}\sigma.\pi : \forall x{:}\sigma.\varphi}
\qquad
\rul{\Gamma \vdash \pi : \forall x{:}\sigma.\varphi \quad \Gamma \vdash a : \sigma}
    {\Gamma \vdash \pi\,a : \subst{\varphi}{a}{x}}
\\[2.2ex]
\rul{\Gamma, X{:}K \vdash \pi : \varphi}
    {\Gamma \vdash \Lambda X{:}K.\pi : \forall X{:}K.\varphi}
\qquad
\rul{\Gamma \vdash \pi : \forall X{:}K.\varphi \quad \Gamma \vdash P : K}
    {\Gamma \vdash \pi\,P : \subst{\varphi}{P}{X}}
\\[2.2ex]
\rul{\Gamma \vdash \pi_1 : \varphi_1 \quad \Gamma \vdash \pi_2 : \varphi_2}
    {\Gamma \vdash \langle\pi_1,\pi_2\rangle : \varphi_1 \wedge \varphi_2}
\qquad
\rul{\Gamma \vdash \pi : \varphi_1 \wedge \varphi_2}
    {\Gamma \vdash \pi.i : \varphi_i}
\qquad
\rul{\Gamma \vdash \pi : \varphi_i}
    {\Gamma \vdash \iota_i(\pi) : \varphi_1 \vee \varphi_2}
\\[2.2ex]
\rul{\Gamma \vdash \pi : \varphi_1 \vee \varphi_2 \quad
     \Gamma, p_i{:}\varphi_i \vdash \pi_i : \psi \ (i = 1,2)}
    {\Gamma \vdash \mathsf{vcase}(\pi; p_1.\pi_1; p_2.\pi_2) : \psi}
\qquad
\rul{\Gamma \vdash a : \sigma \quad \Gamma \vdash \pi : \subst{\varphi}{a}{x}}
    {\Gamma \vdash \langle a, \pi \rangle : \exists x{:}\sigma.\varphi}
\\[2.2ex]
\rul{\Gamma \vdash \pi : \exists x{:}\sigma.\varphi \quad
     \Gamma, x{:}\sigma, p{:}\varphi \vdash \pi' : \psi \quad
     x, p \notin \fv{\psi}}
    {\Gamma \vdash \mathsf{ecase}(\pi; x\,p.\pi') : \psi}
\qquad
\rul{}{\Gamma \vdash \mathsf{tt} : \top}
\\[2.2ex]
\rul{\Gamma \vdash \pi : \bot \quad \Gamma \vdash \varphi \sprop}
    {\Gamma \vdash \texfalso{\pi}{\varphi} : \varphi}
\qquad
\rul{\Gamma \vdash a : \sigma}
    {\Gamma \vdash \mathsf{refl}_{\sigma}(a) : a =_{\sigma} a}
\\[2.2ex]
\rul{\Gamma \vdash \pi : a =_{\sigma} b \quad
     \Gamma, z{:}\sigma \vdash \psi \sprop \quad
     \Gamma \vdash \pi' : \subst{\psi}{a}{z}}
    {\Gamma \vdash \mathsf{rew}(\pi; z.\psi; \pi') : \subst{\psi}{b}{z}}
\end{array}
\]
\caption{Typing of proofs, part I.}
\label{fig:proofs}
\end{figure}

\begin{figure}[tp]
\small
\[
\begin{array}{c}
\rul{\Gamma \vdash \pi :
       \wcns{c}(\tup{a};\tup{\pi}_1) =_{\sigma} \wcns{d}(\tup{b};\tup{\pi}_2)
     \quad \wcns{c} \neq \wcns{d} \text{ constructors of a datatype } I}
    {\Gamma \vdash \mathsf{discr}(\pi) : \bot}
\\[2.2ex]
\rul{\Gamma \vdash t : \refty{x{:}\sigma}{\varphi}}
    {\Gamma \vdash \tvalid{t} : \subst{\varphi}{\tval{t}}{x}}
\qquad
\rul{\Gamma \vdash \pi : \varphi \quad \varphi \conv \varphi' \quad
     \Gamma \vdash \varphi' \sprop}
    {\Gamma \vdash \pi : \varphi'}
\\[2.2ex]
\rul{\begin{array}{c}
     \wcns{k} : \forall \tup{x}{:}\tup{\tau}.\,\varphi_1 \to \dots \to
       \varphi_\ell \to J(\tup{u}) \in \Env \\
     \Gamma \vdash a_i : \subst{\tau_i}{\tup{a}_{<i}}{\tup{x}_{<i}} \
       (\text{all } i) \quad
     \Gamma \vdash \pi_l : \subst{\varphi_l}{\tup{a}}{\tup{x}} \ (\text{all } l)
     \end{array}}
    {\Gamma \vdash \wcns{k}(\tup{a}; \tup{\pi}) :
       J(\subst{\tup{u}}{\tup{a}}{\tup{x}})}
\\[2.2ex]
\rul{\Gamma \vdash \pi : J(\tup{a}) \quad
     \Gamma, \tup{z}{:}\tup{\sigma} \vdash \psi \sprop \quad
     \Gamma \vdash \pi_j : \Xi_j(\tup{z}.\psi) \ \ (\text{all clauses } j)}
    {\Gamma \vdash \tind{J}(\tup{z}.\psi; \tup{\pi}; \pi) :
       \subst{\psi}{\tup{a}}{\tup{z}}}
\\[2.2ex]
\rul{\begin{array}{c}
     \Gamma \vdash a : I(\tup{u}) \qquad
     \Gamma, \tup{z}{:}\tup{\sigma}, z{:}I(\tup{z}) \vdash \psi \sprop \\
     \Gamma,\ \tup{y}_i{:}\tup{A}_i,\ \tup{p}_i{:}\tup{\varphi}_i,\
       \tup{h}_i{:}\bigl(\psi[\tup{v}_{ij}/\tup{z},
       y_{ij}/z]\bigr)_{j \in \mathrm{rec}(i)} \\
       \vdash\ \pi_i :
       \psi[\tup{t}_i/\tup{z},\, \wcns{c}_i(\tup{y}_i;\tup{p}_i)/z]
       \ \ (\text{all } i)
     \end{array}}
    {\Gamma \vdash \tind{I}(\tup{z}.z.\psi;\,
       (\tup{y}_i\,\tup{p}_i\,\tup{h}_i.\pi_i)_i;\, a) :
       \psi[\tup{u}/\tup{z},\, a/z]}
\end{array}
\]
\caption{Typing of proofs, part II.  For a predicate clause
  $\wcns{k}_j : \forall \tup{x}{:}\tup{\tau}.\varphi_1 \to \dots \to
  \varphi_\ell \to J(\tup{u})$, the \emph{clause obligation} is
  $\Xi_j(\tup{z}.\psi) = \forall \tup{x}{:}\tup{\tau}.\,\varphi_1 \to \dots
  \to \varphi_\ell \to \subst{\psi}{\tup{v}_1}{\tup{z}} \to \dots \to
  \subst{\psi}{\tup{v}_s}{\tup{z}} \to \subst{\psi}{\tup{u}}{\tup{z}}$,
  where $J(\tup{v}_1), \dots, J(\tup{v}_s)$ are the recursive premises
  among the $\varphi_i$.  In the datatype induction $\tind{I}$,
  $\mathrm{rec}(i)$ is the set of recursive argument positions of
  $\wcns{c}_i$, and for $j \in \mathrm{rec}(i)$ with $A_{ij} =
  I(\tup{v}_{ij})$ the hypothesis $h_{ij}$ has the displayed type; the
  clause proofs bind $\tup{y}_i, \tup{p}_i, \tup{h}_i$ because the
  conclusion instance mentions the constructor term
  $\wcns{c}_i(\tup{y}_i;\tup{p}_i)$, which requires names for the premise
  proofs.}
\label{fig:proofs2}
\end{figure}

\begin{definition}[Forced singletons]
\label{def:forced-singleton}
An inductive predicate $J : (\tup{\sigma}) \to \Prop$ is a \emph{forced
singleton} if it has exactly one clause
\[
  \wcns{k} : \forall \tup{x}{:}\tup{\tau}.\
  \varphi_1 \to \dots \to \varphi_\ell \to J(\tup{u}^0)
\]
and there is a function $\jmath : \{1, \dots, |\tup{x}|\} \to \{1, \dots,
m\}$, called a \emph{solution}, with $u^0_{\jmath(i)} = x_i$ (syntactic
identity) for every $i$.  Note that $\jmath$ is automatically injective
($u^0_{\jmath(i)}$ determines $x_i$), that repeated variables among other
index positions are permitted (non-linear index patterns), and that the
remaining index positions may hold arbitrary terms.  For an occurrence
$J(\tup{u})$ the induced \emph{solution substitution} is $\theta =
[u_{\jmath(1)}/x_1, \dots, u_{\jmath(|\tup{x}|)}/x_{|\tup{x}|}]$.  In the
$\tscase{J}{}{}{}$ rule of \cref{fig:terms2}, $\theta(\tup{\varphi})$ and
$\theta(\tup{u}^0)$ denote the simultaneous application of $\theta$.
When the elimination is used we consider $\jmath$ fixed once per
declaration (any choice works; \cref{rem:solution-choice}).
\end{definition}

\begin{example}
$\mathsf{eq}_\sigma$ is a forced singleton: the clause
$\wcns{refl}_{\mathsf{eq}} : \forall x{:}\sigma.\ \mathsf{eq}_\sigma(x,x)$
has $\tup{u}^0 = (x, x)$ and $\jmath(1) = 1$ (or $2$).  For an occurrence
$\mathsf{eq}_\sigma(a,b)$, the solution substitution is $\theta = [a/x]$,
so $\tscase{\mathsf{eq}_\sigma}{\pi}{z_1 z_2.\tau}{.t}$ requires $t :
\tau[a/z_1, a/z_2]$ and produces $\tau[a/z_1, b/z_2]$ --- precisely
transport in the second index.  The primitive $\tcast{}{}{}$ is the same
elimination for the primitive equality.  The predicate $\mathsf{le}$ is
not a forced singleton (two clauses); $\mathsf{Acc}$-style predicates are
excluded from $\CICi$ altogether because their recursive premises sit
under quantifiers (\cref{sec:calculus-cic}).
\end{example}

\begin{remark}[Choice of solution]
\label{rem:solution-choice}
If $\jmath$ and $\jmath'$ are both solutions, the two instantiations of
the $\tscase{}{}{}{}$ rule are interderivable: from $\pi : J(\tup{u})$
one obtains, by $\tind{J}$ with motive $\tup{z}.\,(z_{\jmath(i)}
=_{\sigma_{\jmath(i)}} z_{\jmath'(i)})$ --- whose clause obligation is
$u^0_{\jmath(i)} = u^0_{\jmath'(i)}$, i.e.\ $x_i = x_i$, discharged by
$\mathsf{refl}$ --- a proof of $u_{\jmath(i)} = u_{\jmath'(i)}$ for
every $i$, and $\tcast{}{}{}$ along these equations mediates between
the two premise and motive instantiations.  (The equation is well
formed because both index types are convertible to the type of $x_i$.)
Semantically both instantiations are validated by the same argument
(\cref{lem:fundamental}), so fixing $\jmath$ loses no generality.
\end{remark}

\begin{remark}[Relation to CIC's singleton criterion]
\label{rem:singleton-criterion}
CIC permits a match on a $\Prop$-sorted scrutinee in an informative
position when the inductive is empty or a \emph{singleton}: one
constructor whose arguments are all propositional, \emph{parameters put
aside} \cite{Rocq,Letouzey2003}.  Under the monomorphic stratification of
\cref{sec:calculus-cic}, the parameters of an equality-like family become
clause variables that occur verbatim among the result indices, and the
remaining CIC constructor arguments become the premises
$\tup{\varphi}$.  Every CIC singleton elimination is therefore an
instance of $\tscase{}{}{}{}$.  A taxonomy note: what
\cref{def:forced-singleton} demands is that the clause variables be
\emph{verbatim-forced} --- literally among the indices.  The forcing
analyses of Brady et al.\ \cite{BradyEtAl2004} and Gilbert et al.\
\cite{GilbertEtAl2019} are more liberal: they also count
\emph{pattern-forced} arguments, occurring under constructors in an
index term ($x$ under $\csS$, recoverable by injectivity).
$\tscase{}{}{}{}$ therefore sits strictly between CIC's parameter
criterion and the Gilbert et al.\ criterion.  The pattern-forced
generalisation appears semantically available --- index reflection
(\cref{lem:index-reflection}) applied to the index datatype recovers
the argument --- but the solution substitution $\theta$ would no
longer be a syntactic substitution, and we leave it as future work.
In the other direction, the converse of the inclusion above fails:
\cref{def:forced-singleton} does not require the forced variables
to occur \emph{uniformly} in recursive premises, so it accepts clauses
that CIC's parameter mechanism cannot express (e.g.\ $\wcns{k} : \forall
x{:}\nat.\ J(\csS\,x, \csO) \to J(x, \csS\,\csO)$).  The permissiveness
is deliberate and harmless: soundness is proved against the model
(\cref{lem:fundamental}), which validates the general rule; only proof
\emph{reconstruction} of uses of the extra generality requires
inversion-style reasoning rather than a primitive CIC match
(\cref{sec:limitations-reconstruction}).  We note that the intended
front end never \emph{produces} that extra generality: it elaborates
CIC singleton eliminations, whose images satisfy the parameter
criterion.  The margin is robustness of the calculus, not a feature
the pipeline exercises, so the corresponding reconstruction
obligation is expected to be vacuous in practice.  Empty elimination needs no new
former: for a predicate $J$ declared with zero clauses, $\tind{J}$ with
motive $\bot$ and no clause obligations yields $\bot$ from $J(\tup{a})$,
and $\texfalso{}{}$ concludes.
\end{remark}

\subsection{Reduction and conversion}
\label{sec:calculus-red}

\begin{definition}[Reduction]
\label{def:source-red}
The relation $\red$ on terms is the contextual closure --- in all term,
proof, type, proposition and predicate positions --- of the following
root rules:
\begin{align*}
  (\lambda x{:}\sigma.t)\,a &\red \subst{t}{a}{x}
  & (\lambda p{:}\varphi.t)\,\pi &\red \subst{t}{\pi}{p}
  \\
  \tfix\,f{:}T.\,t &\red \subst{t}{\tfix\,f{:}T.\,t}{f}
  & \tval{\trefine{t}{\pi}} &\red t
  \\
  c &\red t \quad \text{if } (c \defeq t : \sigma) \in \Env
  & \tcast{\mathsf{refl}_\sigma(a)}{z.\tau}{t} &\red t
  \\
  \tsplit{\langle \pi_1, \pi_2 \rangle}{p_1\,p_2.t} &\red
     \subst{\subst{t}{\pi_1}{p_1}}{\pi_2}{p_2}
\end{align*}
together with the singleton-elimination rule
\[
  \tscase{J}{\wcns{k}(\tup{a};\tup{\pi})}{\tup{z}.\tau}{\tup{p}.t}
     \red \subst{t}{\tup{\pi}}{\tup{p}}
\]
and the rule for case analysis meeting a constructor,
\begin{multline*}
  \tcase{\wcns{c}_i(\tup{a};\tup{\pi})}{\tup{z}.z.\tau}{\{\dots\}}
     \red\\
     t_i[\tup{a}/\tup{y}_i,\ \tup{\pi}/\tup{p}_i,\
        \overrightarrow{\mathsf{refl}}(\tup{t}_i[\tup{a}/\tup{y}_i])/\tup{q}_i],
\end{multline*}
where $\overrightarrow{\mathsf{refl}}(\tup{t}_i[\tup{a}/\tup{y}_i])$
abbreviates the proof list
\[
  \mathsf{refl}_{\sigma_1}(t_{i1}[\tup{a}/\tup{y}_i]),\ \dots,\
  \mathsf{refl}_{\sigma_m}(t_{im}[\tup{a}/\tup{y}_i])
\]
substituted for the index-equation variables.  Source \emph{conversion}
$\conv$ is the least
congruence (on each syntactic category, closing under all formers)
containing $\red$.
\end{definition}

Reduction is untyped and unrestricted.  As in the companion paper, no
metatheoretic property of $\red$ itself is needed anywhere: $\conv$
enters only through the conversion typing rules, and the semantics of
\cref{sec:semantics} is invariant under it (\cref{lem:conv-inv}).  Note
that the index-equation variables $\tup{q}_i$ receive
\emph{reflexivity} proofs when a case meets a constructor: at the redex,
the occurrence indices are convertible to the constructor's index
pattern, and the conversion typing rule absorbs the mismatch.  Since
proofs are erased (\cref{def:erasure}), this substitution is invisible
in the target.

\begin{lemma}[Rigidity of type formers]
\label{lem:rigid}
If $\sigma \conv \sigma'$ then $\sigma$ and $\sigma'$ have the same
outermost former: both are occurrences $I(\tup{u})$, $I(\tup{u}')$ of the
same datatype $I$ with $u_h \conv u_h'$ for all $h$, or both are $\Pi
x{:}\rho.\tau$-types whose domains and codomains are respectively
convertible, or both are $\varphi \to \tau$-types with convertible
components, or both are refinement types with convertible components.
The analogous statement holds for propositions.
\end{lemma}

\begin{proof}
No root rule of \cref{def:source-red} applies at the root of a set type
or a proposition --- the root rules rewrite terms and proofs only --- so
every $\red$-step inside a type or proposition happens strictly below its
outermost former: inside an index term of a datatype occurrence, inside a
component type or proposition, or inside an embedded term or proof.  The
claim follows by induction on the length of the conversion path,
observing that a single step preserves the outermost former and relates
components (index terms included) pointwise.
\end{proof}

\begin{lemma}[Substitution]
\label{lem:source-subst}
All judgements admit substitution: if $\Gamma, x{:}\sigma, \Gamma'
\vdash \mathcal{J}$ and $\Gamma \vdash a : \sigma$, then $\Gamma,
\subst{\Gamma'}{a}{x} \vdash \subst{\mathcal{J}}{a}{x}$, and analogously
for proof variables and predicate variables.
\end{lemma}

\begin{proof}
By induction on the derivation of the second judgement, simultaneously
for all judgement forms and all three substitution kinds; write
$[\cdot]$ for the substitution being pushed ($[a/x]$ in the notation
of the statement).  Three mechanisms carry the routine cases, and two
new rules deserve their arguments in full.

\emph{Routine mechanisms.}  (i) Every rule is homomorphic in its
components, so the induction hypotheses on the premises reassemble the
rule instance, using the composition identity
$\tau[a'/y][a/x] = \subst{\tau}{a}{x}[\,\subst{a'}{a}{x}/y]$ (for $y
\neq x$, $y \notin \fv{a}$, after $\alpha$-renaming) wherever the
conclusion contains a substituted phrase; in the constructor, case and
$\tind{I}$ rules the declaration data ($A_{ij}$, $\varphi_{il}$,
$\tup{t}_i$) mention only the constructor telescope and earlier
declarations, so they are untouched by $[\cdot]$ and only the
instantiating terms absorb it, e.g.\
$(A_{ij}[\tup{a}_{<j}/\tup{y}_{<j}])[a/x] =
A_{ij}[(\tup{a}_{<j}[a/x])/\tup{y}_{<j}]$.  (ii) $\conv$ is closed
under substitution: each root rule of \cref{def:source-red} is ---
for the forded-case rule, both sides substitute componentwise and the
reflexivity list $\overrightarrow{\mathsf{refl}}
(\tup{t}_i[\tup{a}/\tup{y}_i])$ becomes the reflexivity list of the
substituted instances --- so the two conversion rules follow from the
induction hypothesis.  (iii) Substitution of a predicate expression
into an atom, $\subst{X(\tup{a})}{P}{X} = P \cdot \tup{a}$ (predicate
application, $\beta$-reducing the $\lambda$-form), yields a well-formed
proposition by the kind premise, and the process terminates because
$P$ cannot mention $X$ after $\alpha$-renaming.

\emph{The $\tscase{}{}{}{}$ case (commutation of $\theta$ with the
ambient substitution).}  Let the rule be applied at an occurrence
$J(\tup{u})$, with clause data $\wcns{k} : \forall
\tup{x}{:}\tup{\tau}.\ \varphi_1 \to \dots \to \varphi_\ell \to
J(\tup{u}^0)$, solution $\jmath$, and solution substitution $\theta =
[u_{\jmath(i)}/x_i]_i$.  After the ambient substitution the scrutinee
proof has type $J(\tup{u}[a/x])$, whose solution substitution is
$\theta' = [\,\subst{u_{\jmath(i)}}{a}{x}/x_i]_i$ (the clause data
$\tup{u}^0, \jmath$ are closed declaration components and do not
change).  For every clause phrase $\xi$ --- a premise $\varphi_l$ or
the pattern tuple $\tup{u}^0$ --- the free variables of $\xi$ lie
among $\tup{x}$ (\cref{def:env}), so the composed substitution acts on
$\xi$ only through the images of $\tup{x}$:
\[
  (\theta(\xi))[a/x] \;=\; \xi\bigl[\subst{u_{\jmath(i)}}{a}{x}/x_i
  \bigr]_i \;=\; \theta'(\xi).
\]
Hence $\theta(\tup{\varphi})[a/x] = \theta'(\tup{\varphi})$ and
$\tau[\theta(\tup{u}^0)/\tup{z}][a/x] =
\subst{\tau}{a}{x}[\theta'(\tup{u}^0)/\tup{z}]$ (using (i) for the
motive), and the induction hypotheses re-derive the rule at
$J(\tup{u}[a/x])$ with solution substitution $\theta'$, producing the
substituted conclusion type $\subst{\tau}{a}{x}[\tup{u}[a/x]/\tup{z}]
= (\tau[\tup{u}/\tup{z}])[a/x]$.

\emph{The $\mathsf{discr}$ case.}  From $\Gamma, x{:}\sigma, \Gamma'
\vdash \pi : \wcns{c}(\tup{a}';\tup{\pi}_1) =_\sigma
\wcns{d}(\tup{b}';\tup{\pi}_2)$ with $\wcns{c} \neq \wcns{d}$:
substitution is homomorphic on constructor terms,
$(\wcns{c}(\tup{a}';\tup{\pi}_1))[a/x] = \wcns{c}(\tup{a}'[a/x];
\tup{\pi}_1[a/x])$, so both equation sides remain constructor-headed
with the same distinct head constructors, and the rule re-applies to
the substituted premise supplied by the induction hypothesis.  (This
is the only rule whose side condition inspects the \emph{shape} of a
term, which is why it is singled out.)
\end{proof}

\subsection{The shallow fragment}
\label{sec:calculus-shallow}

\begin{definition}[Shallow fragment $\Sfrag$]
\label{def:shallow-fragment}
A proposition is \emph{shallow} if it contains no subformula $\forall
X{:}K.\varphi$ and no atom $X(\tup{a})$ with $X$ a predicate variable; a
set type is shallow if all propositions occurring in it (refinement
payloads and premises) are shallow.  $\Sfrag$ is the class of shallow
propositions and set types.
\end{definition}

Index terms of datatype occurrences are terms and are always shallow
material; only genuine predicate quantification is excluded.  As in the
companion paper, the translation is defined on $\Sfrag$; proofs inside
the environment may use the full calculus, but a lemma whose statement
leaves $\Sfrag$ contributes no premise axiom, a definition whose type
leaves $\Sfrag$ contributes no specification axiom, and goals must be
shallow and closed.  We require the declarations of $\Env$ (datatype
index types, constructor argument types, premises and index terms;
inductive-predicate components) to lie in $\Sfrag$.

\subsection{Programming with forded case analysis}
\label{sec:calculus-programming}

The index equations $\tup{q}_i$ in branches are what makes dependently
typed idioms expressible with \emph{first-order} motives.  We record the
two standard examples and the derivability of injectivity; both are used
in \cref{sec:translation} as running translation examples.

\begin{lemma}[Derivable injectivity]
\label{lem:inj-derivable}
Let $\wcns{c}$ be a constructor of a datatype $I$ whose $j$-th argument
type $A_j$ is closed --- it mentions no earlier telescope variables; it
may be a non-recursive type or a recursive occurrence $I(\tup{v})$ at
closed indices (for a non-indexed datatype, in particular, any
recursive position qualifies).  Then from $\Gamma \vdash \pi :
\wcns{c}(\tup{a};\tup{\pi}_1) =_\sigma \wcns{c}(\tup{b};\tup{\pi}_2)$
one can derive $\Gamma \vdash \mathsf{inj}_j(\pi) : a_j =_{A_j} b_j$
for a proof term $\mathsf{inj}_j(\pi)$ definable uniformly from the
derivation.
\end{lemma}

\begin{proof}
Since $\Gamma \vdash \wcns{c}(\tup{a};\tup{\pi}_1) : \sigma$ and the
constructor rule produces the type $I(\tup{t}[\tup{a}/\tup{y}])$, by
\cref{lem:rigid} $\sigma \conv I(\tup{u})$ for some $\tup{u}$; fix such
an occurrence.  Define the projection
\[
  F \defeq \lambda w{:}I(\tup{u}).\ \tcase{w}{\tup{z}.z.A_j}
    {\{\wcns{c}(\tup{y};\tup{p};\tup{q}) \Rightarrow y_j \mid
      \wcns{d}(\dots) \Rightarrow a_j \ (\wcns{d} \neq \wcns{c})\}}.
\]
The motive is the constant type $A_j$, well formed in the motive
context since $A_j$ is closed (this includes the recursive case $A_j =
I(\tup{v})$ with $\tup{v}$ closed: a legal occurrence like any other);
the $\wcns{c}$-branch is well typed because $y_j : A_j$ (again by
closedness, no instantiation of earlier telescope variables occurs) and
the branch obligation is $A_j$ itself; the default branches use $a_j$,
well typed over $\Gamma$ by weakening (its type
$A_j[\tup{a}_{<j}/\tup{y}_{<j}] = A_j$ by closedness).  Then $F$ is
well typed, $F\,\wcns{c}(\tup{a};\tup{\pi}_1) \conv a_j$ and
$F\,\wcns{c}(\tup{b};\tup{\pi}_2) \conv b_j$ by a $\beta$- and a
case-step.  Now
\[
  \mathsf{inj}_j(\pi) \defeq
  \mathsf{rew}\bigl(\pi;\ z.\,(a_j =_{A_j} F\,z);\
    \mathsf{refl}_{A_j}(a_j)\bigr),
\]
where the base proof is well typed at $(a_j =_{A_j}
F\,\wcns{c}(\tup{a};\tup{\pi}_1)) \conv (a_j =_{A_j} a_j)$ by the
conversion rule, and the conclusion converts to $a_j =_{A_j} b_j$.
\end{proof}

\begin{example}[Head and tail]
\label{ex:vhead-vtail}
With the declarations of \cref{ex:declarations}:
\begin{align*}
  \mathsf{vhead} \defeq{} & \lambda n{:}\nat.\ \lambda
    v{:}\vect(\csS\,n).\\
  &\tcase{v}{\tup{z}.z.A}{\{
      \csvnil(;;q) \Rightarrow \texfalso{\mathsf{discr}(q)}{A} \mid
      \csvcons(k, a, w;; q) \Rightarrow a\}}\\[0.7ex]
  \mathsf{vtail} \defeq{} & \lambda n{:}\nat.\ \lambda
    v{:}\vect(\csS\,n).\\
  &\tcase{v}{\tup{z}.z.\vect(n)}{\bigl\{
    \begin{aligned}[t]
      &\csvnil(;;q) \Rightarrow \texfalso{\mathsf{discr}(q)}{\vect(n)}
        \mid\\
      &\csvcons(k, a, w;; q) \Rightarrow{}\\
      &\qquad
        \tcast{\mathsf{sym}(\mathsf{inj}_1(q))}{z.\vect(z)}{w}\
        \bigr\}
    \end{aligned}}
\end{align*}
where $\mathsf{sym}(\pi) \defeq \mathsf{rew}(\pi;\, z.\,(z =_\nat n);\,
\mathsf{refl}_\nat(n))$ derives $k =_\nat n$ from $\mathsf{inj}_1(q) : n
=_\nat k$, and the $\csvcons$ branch of $\mathsf{vtail}$ is, written in
full, $\tcast{\mathsf{sym}(\mathsf{inj}_1(q))}{z.\vect(z)}{w}$.  In the
$\csvnil$ branches, $q : \csS\,n =_{\nat} \csO$ and $\mathsf{discr}(q) :
\bot$ discharges the impossible branch --- with a \emph{constant}
motive; no large elimination, no $\mathsf{unit}$-valued return-type
trick.  In the $\csvcons$ branch of $\mathsf{vtail}$, $q : \csS\,n
=_\nat \csS\,k$ gives $\mathsf{inj}_1(q) : n =_\nat k$ by
\cref{lem:inj-derivable}, whence $\mathsf{sym}(\mathsf{inj}_1(q)) : k
=_\nat n$, and the cast transports $w : \vect(k)$ to $\vect(n)$.  Both
definitions are well typed in $\CICi$; their translations are computed
in \cref{ex:vhead-transl}.
\end{example}

\subsection{Relation to CIC}
\label{sec:calculus-cic}

$\CICi$ is designed as the image of a first-order, monomorphic slice of
CIC under analyses that are part of the CoqHammer translation pipeline
or of program extraction.  The discussion of stratification,
monomorphisation, refinement classification, structural fixpoints and
the propositional connectives in \cite[\S 2.6]{Extraction2026} applies
verbatim; we record only what is new for indexed families.

\begin{itemize}[leftmargin=2em]
\item \emph{Parameters and indices.}  CIC distinguishes uniform
  parameters from indices; the distinction matters for elimination (a
  parameter is fixed across a match, an index is generalized).  After
  monomorphisation, \emph{type} parameters are ground and disappear into
  the datatype name ($\vect$ above is really
  $\mathsf{vec}_A$).  \emph{Term} and \emph{proposition} parameters
  become indices ($\mathsf{eq}_\sigma$'s left argument) or premises
  ($\mathsf{rfl}_{\varphi_0}$'s $P$), and the forced-variable analysis of
  \cref{def:forced-singleton} recovers, per occurrence, exactly the
  specialization that CIC's parameter mechanism provided --- this is the
  content of \cref{rem:singleton-criterion}.  Non-uniform parameters are
  indices here, as they are in CIC's kernel.
\item \emph{Dependent case analysis.}  CIC's
  $\mathsf{match}\ a\ \mathsf{in}\ I\ \tup{z}\ \mathsf{return}\
  \tau\ \mathsf{with}\ \dots$ has branches typed at each constructor's
  own indices, with no equational premises; impossible branches are
  discharged by large-eliminating motives.  The forded case of $\CICi$
  instead hands each branch the index equations.  The two styles are
  intertranslatable in CIC extended with UIP --- this is precisely the
  Goguen--McBride--McKinna elaboration \cite{GoguenEtAl2006} --- and the
  front end is expected to produce forded branches by the same analysis
  Coq's \texttt{inversion} performs \cite{CornesTerrasse1995,Monin2010}:
  generalize the indices, introduce equations, and (optionally) simplify
  them.  Branches the front end can discharge by rigid-index conflict
  arrive as $\texfalso{\mathsf{discr}(\cdot)}{\cdot}$ as in
  \cref{ex:vhead-vtail}; branches it cannot discharge (green slime)
  arrive with their equations intact, which is the point
  (\cref{rem:slime}).  Large elimination remains excluded from $\CICi$:
  fording removes its main \emph{programming} use (impossible branches),
  and its \emph{logical} uses (proving discrimination) are covered by
  the primitive $\mathsf{discr}$, which is CIC-derivable by exactly such
  a large elimination.
\item \emph{Propositional index positions.}  Index types in $\CICi$ are
  set types; CIC families indexed by \emph{propositions} or \emph{proofs}
  have no direct image.  The front end drops proof-sorted index
  positions before entering the fragment; this is sound for the same
  reason the model is proof-irrelevant (all proofs of a proposition
  denote one element --- \cref{sec:semantics}), and is the
  translation-level counterpart of the rule that erased data may appear
  in indices only where nothing relevant observes them
  \cite{Atkey2018,GilbertEtAl2019}.  Proof \emph{subterms} of index
  terms, by contrast, are in the fragment and are erased structurally
  (\cref{rem:proof-indices}).
\item \emph{Recursion.}  The $\tfix$ former captures structural
  recursion on a family element with a designated guard argument.
  Recursive calls may be at arbitrary index instantiations (the call
  $f\,k\,w$ inside $\mathsf{vmap}$-style definitions instantiates the
  index argument with the pattern variable $k$); the guard condition
  constrains only the guard position.  Well-founded
  ($\mathsf{Acc}$-based) recursion remains excluded, for the reasons
  analysed in \cite[\S 10.3]{Extraction2026}: its unfolding equations
  are sound only with their propositional premises retained, and its
  guardedness is semantic rather than syntactic
  \cite{GilbertEtAl2019}.  Relatedly, predicate clauses with recursive
  premises under quantifiers ($\mathsf{Acc}$'s constructor) are excluded
  from the declaration format.
\end{itemize}
